\chapter*{Conclusion Générale}
\addcontentsline{toc}{chapter}{Conclusion Générale}

Ce travail a abouti à \appname, une application de bureau de gestion académique conçue pour fonctionner
intégralement sur le réseau local d'un établissement, sans connexion Internet. Le bilan peut être lu à
trois niveaux.

Sur le plan fonctionnel, l'application couvre la chaîne complète de la scolarité : inscription et dossier
de l'élève, structure pédagogique et calendrier, emploi du temps et volumes horaires, séances,
présences, observations, récompenses et sanctions, examens, surveillance, notes et moyennes, supports de
cours, bibliothèque documentaire avec revue éditoriale, messagerie interne, notifications et rapports
imprimés. Les sept rôles identifiés disposent
chacun d'un périmètre cohérent, et les opérations sensibles sont contrôlées par le serveur plutôt que par
l'écran qui les affiche.

Sur le plan technique, la séparation entre PostgreSQL, le processus principal et l'interface React
explique les propriétés obtenues : une donnée unique partagée par tous les postes, des règles
d'écriture centralisées, une interface qui reste réactive et une génération de documents qui ne dépend
d'aucun service externe. L'internationalisation à trois langues, avec une lecture de droite à gauche
réellement prise en charge, des libellés de base traduits et des dates qui suivent la langue de
l'interface, répond à une exigence courante des établissements et aurait été difficile à ajouter après
coup.

Sur le plan de l'exploitation, le déploiement a été pensé pour un technicien et non pour un
développeur : un préréglage machine évite la saisie de la connexion sur chaque poste, un écran de
configuration précède la connexion, les fichiers partagés sont référencés par des clés relatives, les
mises à jour se font depuis le partage local avec vérification par empreinte, et la sauvegarde de la
base est scriptée.

Le travail comporte néanmoins des limites qu'il faut nommer. La première est l'absence de suite de
tests automatisés dans le dépôt : la validation repose sur des scénarios manuels documentés dans le
guide de déploiement, ce qui rend une régression plus coûteuse à détecter qu'elle ne le serait avec des
tests de services et de rendu. La deuxième est la dépendance à un serveur unique : si PostgreSQL ou le
partage de fichiers s'arrête, tous les postes perdent l'accès aux données et aux médias, même si le
client continue de fonctionner et signale l'incident. La troisième concerne l'internationalisation :
les messages produits par le processus principal, construits à partir de la logique métier et des
erreurs de base, restent en anglais et ne passent pas par les dictionnaires de l'interface. La
quatrième est l'ergonomie de certaines pages très denses, qui suppose une courte formation pour les
utilisateurs peu familiers des tableaux de gestion.

Ces limites n'invalident pas l'apport principal : l'établissement dispose désormais d'une source unique
pour les dossiers, les séances et les résultats, consultable en temps réel par les services concernés,
sur un réseau qui ne sort pas de ses murs.
